\section{Classical feedback and coherent reset width}
\label{sec:reset88}
The independent blocks of Definition~\ref{def:blocks87} do not
exhaust the strategies with bounded coherent acquisition.  We now
allow both classical feedback between blocks and adaptive quantum
control within a block.  The receiver may retain and process quantum
memory throughout.  The restriction is on importing that memory
into the next acquisition, not on its dimension or final readout.

\begin{definition}[Reset protocols]\label{def:reset88}
Fix integers $1\le b\le N$.  At a classical history $h$, a protocol
may stop, or choose a positive integer $n(h)\le b$ and prepare a new
complete probe--reference system in a state $\sigma_h$.  Conditional
on $h$, this state is the same under both hypotheses and is in a
tensor product with the receiver's entire existing memory.  A common
$n(h)$-call adaptive experiment acts on the new system alone; its
internal inputs, quantum memory and controls are unrestricted.
After acquisition, its whole accessible output is delivered to the
receiver.  A common quantum instrument may act jointly on this
output and all previous receiver memory, retaining quantum memory
and producing the next classical history.  No quantum subsystem of
that memory enters a subsequent acquisition.  On every history,
\[
 \sum_r n(h_r)\le N.
\]
The integers are reserved call budgets: unused internal slots can
be padded and are still charged.  Public stopping at block boundaries
is allowed.  All classical histories are finite strings with finite
or countable outcome sets; conditional state preparations and
instruments are defined also at histories null under one hypothesis.
Let $D_N^{\mathrm{reset},b}(E,F)$ be the supremum of the unhalved
trace separation of the final accessible states, including the
classical record and arbitrary final joint processing.
\end{definition}

The complete new system, including every reference and purification
used in acquisition, must be independent of the receiver conditional
on $h$.  Separability of the probe marginal alone is not sufficient.
The receiver memory need not be classical, and can be entangled with
previous block outputs.  Measuring that memory to choose a fresh
input is allowed; transmitting it coherently into that input is not.
In particular,
\begin{equation}\label{eq:resetinclusions88}
 D_N^{[b]}\le D_N^{\mathrm{reset},b}\le D_N,
 \qquad D_N^{\mathrm{reset},N}=D_N.
\end{equation}
The first inclusion uses no feedback and postpones receiver processing;
the equality uses one unrestricted adaptive block.  At $b=1$ the
receiver can still adapt future preparations through classical
messages.  Equality with the independent-product distance is not
asserted.

For positive, possibly subnormalized operators, write
$f(R,S)=\|\sqrt R\sqrt S\|_1$.  For normalized states put
$d_B(R,S)^2=2(1-f(R,S))$.  We use
\begin{equation}\label{eq:rootrules88}
\begin{split}
 f(R\otimes\rho,S\otimes\sigma)&=f(R,S)f(\rho,\sigma),\\
 f\left(\bigoplus_zR_z,\bigoplus_zS_z\right)&=\sum_z f(R_z,S_z),\\
 f(\mathcal A(R),\mathcal A(S))&\ge f(R,S)
\end{split}
\end{equation}
for a channel $\mathcal A$.  The last assertion for unnormalized
operators follows by separate normalization and homogeneity.
For normalized states,
\begin{equation}\label{eq:tracebures88}
 d_B(R,S)^2\le\|R-S\|_1
 \le2\sqrt{1-f(R,S)^2}\le2d_B(R,S).
\end{equation}
These are standard fidelity identities and inequalities
\cite[Chapter~3]{Watrous65}.

\begin{lemma}[Conditional fidelity accumulation]\label{lem:resetfidelity88}
Consider a reset protocol in which the fresh output states at node
$h$, before receiver processing, satisfy
\[
 f(\omega_{0,h},\omega_{1,h})\ge(1-a_h)_+,
 \qquad a_h\ge0.
\]
Suppose every root-to-terminal history has $\sum_h a_h\le A$.
For any initial receiver operators $R_0,R_1\succeq0$, the final
classical--quantum operators satisfy
\begin{equation}\label{eq:fidelityaccumulation88}
 f(R_{0,\mathrm{out}},R_{1,\mathrm{out}})
 \ge (1-A)_+f(R_0,R_1).
\end{equation}
Thus the receiver's quantum processing and classical feedback need
not preserve a product state across blocks.
\end{lemma}
\begin{proof}
The assertion is immediate for $A\ge1$.  For $A<1$, prove the stronger
continuation statement at each node, for arbitrary incoming positive
operators and a bound $A_h<1$ on the remaining path sum.  At a terminal
node a final common channel only increases fidelity.  At a nonterminal
node let $a=a_h$ and bound every child's remaining sum by $A_h-a$.
Appending the fresh output multiplies the incoming fidelity by at
least $1-a$.  Apply the receiver instrument, retaining its classical
outcome.  If its two subnormalized outputs at outcome $z$ are
$R_{0,z},R_{1,z}$, channel monotonicity and direct-sum additivity give
\[
 \sum_z f(R_{0,z},R_{1,z})
 \ge (1-a)f(R_0,R_1).
\]
The induction hypothesis applies separately at every $z$, even
though the two conditional receiver states differ.  The continued
output fidelity is at least
\[
 [1-(A_h-a)](1-a)f(R_0,R_1)
 \ge(1-A_h)f(R_0,R_1).
\]
Here $0\le a\le A_h<1$, so the product estimate has the stated
sign.  Each acquisition costs at least one call, hence the depth is
at most $N$ and backwards induction terminates.  Countably many
outcomes are handled by the same nonnegative sum; no normalization
at a zero-probability history is required.  Conditional public
randomization is another retained instrument outcome.  This proves
the statement, including stopping and final joint processing.
\end{proof}

\begin{lemma}[A fresh adaptive block]\label{lem:adaptiveblock88}
Fix $E,H,\Lambda$ in Definition~\ref{def:tangenttube86}, and assume
$Q_jH_jQ_j=0$ for every $j$.  Let $A_j=\sqrt{E_j}$ and choose the
canonical factors $B_j$ of Lemma~\ref{lem:curvelogical85} with
$A_j^*B_j+B_j^*A_j=H_j$.  Set
\[
 \beta_0=\|B\|_{\mathrm{op}},\quad
 K=k\beta_0^2(2+\max_j\|H_j\|_{\mathrm{op}}),\quad
 r=\Lambda+K,\quad \alpha=3\beta_0/2.
\]
If $0<s\le1$ and $s\beta_0\le1$, the outputs of every common fresh
$n$-call adaptive experiment obey
\begin{equation}\label{eq:adaptiveblock88}
 d_B(\omega_E,\omega_F)\le s(\alpha n+\sqrt{rn}),\qquad
 1-f(\omega_E,\omega_F)
 \le \tfrac12(\alpha+\sqrt r)^2 n^2s^2,
\end{equation}
for every legal $F$ with $\|F-E-sH\|_\Sigma\le\Lambda s^2$.
The constants are independent of the input, reference, internal
controls and permitted remainder.
\end{lemma}
\begin{proof}
The normalized surrogate in Lemma~\ref{lem:factorremainder86} has
isometries $W_0,W_s$ with $\|W_s-W_0\|\le\alpha s$ and
\[
 \|\mathcal M_F-\mathcal M_{G(s)}\|_\diamond\le rs^2.
\]
Purify the fresh input and all internal common controls.  Inserting
the difference $W_s-W_0$ successively in the $n$ slots bounds the
final purified vector difference by $n\alpha s$: every other factor
is an isometry.  This argument is valid for adaptive controls, not
only tensor-power acquisition.  The discarded dilation environments
are used in the proof but are unavailable to the experiment.
Purification and contraction give the same Bures upper for the
accessible $E$ and surrogate outputs.  Replacing the surrogate by
$F$ costs at most $nrs^2$ in unhalved trace norm by adaptive channel
telescoping.  The first inequality of \eqref{eq:tracebures88} and the
Bures triangle inequality give the first bound in
\eqref{eq:adaptiveblock88}.  The second follows from
$1-f=d_B^2/2$ and $\sqrt n\le n$.  Internal early stopping is padded
within the already reserved $n$ slots and then discarded.  This
retains the full $nrs^2$ error before taking its square root.
\end{proof}

The canonical factor above is generally not horizontal.  Its
linear-in-$n$ purified displacement is appropriate for a coherent
block.  The separate horizontal construction remains essential to
the stronger regular-direction bound for unrestricted adaptive
strategies.

\begin{theorem}[Reset-width interpolation]\label{thm:resetwidth88}
Fix a measurement $E$, a nonzero Hermitian zero-sum tuple $H$ with
$J_j=Q_jH_jQ_j\succeq0$, and $\Lambda\ge0$.  There are
$c,C,s_0>0$, depending only on $E,H,\Lambda$, such that for every
$N\ge1$, $1\le b\le N$, $0<s\le s_0$, and every legal measurement
$F$ satisfying
\[
 \|F-E-sH\|_\Sigma:=\sum_j\|F_j-E_j-sH_j\|_{\mathrm{op}}
 \le\Lambda s^2,
\]
one has $c\mathfrak r\le D_N^{\mathrm{reset},b}(E,F)\le C\mathfrak r$,
where
\begin{equation}\label{eq:resetwidth88}
\mathfrak r=
\begin{cases}
 \min\{1,\sqrt N s\},&H\in\operatorname{Ran}\mathcal C_E,\\
 \min\{1,\sqrt{Nb}\,s\},&J_j=0\ (\forall j),\quad
                         \Gamma_E(H)\notin\mathcal W_E,\\
 \min\{1,Ns\},&J_j\ne0\ (\exists j).
\end{cases}
\end{equation}
Here $\Gamma_E(H)=\frac i2\sum_j[P_j,H_j]$ and
$\mathcal W_E=\sum_jP_j\mathcal H_dP_j$.
All constants are uniform in $N,b,F$, not in $E,H,\Lambda$.
For each $N,b,s$, the lower experiment can be chosen independently
of the unknown remainder.  It can in fact be chosen without
inter-block feedback and with parallel acquisition inside each block.
\end{theorem}
\begin{proof}
In the tangential cases put $\gamma=(\alpha+\sqrt r)^2/2$ as in
Lemma~\ref{lem:adaptiveblock88}.  For a particular public policy let
\[
 Q=\sup_{\text{terminal histories}}\sum_r n(h_r)^2\le bN.
\]
At a node, the fresh experiment has normalized outputs independent
of the incoming receiver memory, conditional on that history.
Lemma~\ref{lem:adaptiveblock88} bounds its fidelity deficit by
$a_h=\gamma n(h)^2s^2$.  Lemma~\ref{lem:resetfidelity88}, starting
from the same normalized initial memory, therefore gives
\[
 f(\rho_{E,\mathrm{out}},\rho_{F,\mathrm{out}})
 \ge(1-\gamma Qs^2)_+.
\]
If $\gamma Qs^2\le1$, use \eqref{eq:tracebures88} and
$1-(1-x)^2\le2x$; otherwise the trace bound two suffices.  In both
cases we obtain the finite policy-sensitive estimate
\begin{equation}\label{eq:resetfinite88}
 D(\text{this protocol};E,F)
 \le\min\{2,\,2(\alpha+\sqrt r)s\sqrt Q\}
 \le\min\{2,\,2(\alpha+\sqrt r)s\sqrt{bN}\}.
\end{equation}
The uniform per-block bound permits every history-dependent choice
of inputs and internal controls, and arbitrary common receiver
instruments.  The two output states are not assumed to be products
across feedback branches.

For the coherent row, the lower construction and the two integer
allocation cases of Theorem~\ref{thm:widthlaw87} belong to the reset
class by \eqref{eq:resetinclusions88}.  Its fixed binary event pays
$m\kappa s^2$ before the finite Bernoulli amplification.  Thus the
same constants and remainder-independent lower apply here.
For the regular row, the product lower and the fully adaptive
horizontal bound \eqref{eq:tubefiniteupper86} sandwich this class.
That upper retains $(\Lambda+K_h)Ns^2$, which is controlled by
$\sqrt N s$ when $\sqrt N s\le1$ and by the trace cap otherwise.
For the opening row, the fixed impossible event
\eqref{eq:leakagelower86} and the adaptive hybrid bound
\eqref{eq:tubehybrid86} give the result already without entanglement.
Take the minimum of the finitely many fixed scale restrictions.
The full range criterion \eqref{eq:fullrange86} makes the cases
disjoint and exhaustive.
\end{proof}

\begin{corollary}[Classical feedback at unit reset width]
\label{cor:classicalfeedback88}
In a fixed coherent tangent neighborhood,
\[
 D_N^{\mathrm{reset},1}(E,F)\asymp\min\{1,\sqrt N s\},\qquad
 D_N^{\mathrm{par}}(E,F)\asymp D_N(E,F)\asymp\min\{1,Ns\}.
\]
The first class permits arbitrary receiver quantum memory and
classical feedback, but no coherent import of that memory into a
fresh probe--reference preparation.  These are local comparison
orders, not equality of exact distances or fixed-pair error exponents.
\end{corollary}
\begin{proof}
Apply Theorem~\ref{thm:resetwidth88} at $b=1$ and combine the
parallel endpoint in Corollary~\ref{cor:parallel87} with the adaptive
tube theorem.  All intervals and constants refer to the same fixed
$E,H,\Lambda$.
\end{proof}

The role of feedback in channel discrimination has extensive prior
study \cite[Theorem~9 and Remark~10]{SHW87}.  Our assertion is a
finite, remainder-uniform consequence of a conditional fidelity
estimate, with arbitrary receiver memory and variable reserved block
sizes.  It does not introduce the classical-feedback resource model
or the standard-versus-coherent metrological mechanism.  The
isometric-extension estimates likewise have direct antecedents in
channel Bures bounds \cite[Theorems~17 and~19]{HMNW88}.

\subsection{Finite policy certificates}
For a finite public decision tree, its pathwise call maximum and
quadratic acquisition cost are computed recursively.  At a terminal
node set $N_h=Q_h=0$; at an acquisition node set
\[
 N_h=n(h)+\max_zN_{hz},\qquad
 Q_h=n(h)^2+\max_zQ_{hz}.
\]
These are maxima over complete histories, not expected costs under
either hypothesis.  Exact integer arithmetic verifies $N_h\le N$
and $n(h)\le b$ on the whole policy.  Given rational $s$ and a
separately established rational majorant $\bar\gamma\ge\gamma$,
\[
 f_{\rm out}\ge(1-\bar\gamma Qs^2)_+,
 \qquad \|\rho_{E,\rm out}-\rho_{F,\rm out}\|_1^2
       \le\min\{4,8\bar\gamma Qs^2\}.
\]
The same induction also gives the stronger recursion
$F_h=(1-\bar\gamma n(h)^2s^2)_+\min_z F_{hz}$ with terminal
value one.  A policy certificate checks these budget consequences.  It does not
verify the physical reset condition or establish the supplied
one-block fidelity majorant.  Those remain hypotheses.  In
particular a numerical or declared reference dimension is not a
certificate of conditional independence.

The explicit nonempty interval $s_*$ of
Proposition~\ref{prop:nonempty87} continues to certify legality and
the sufficient allowance only.  It is not the discrimination
interval $s_0$ above; that interval also depends on the lower witness
and logical correction gap.  The sufficient $\Lambda_*$ is neither
minimal nor canonical.  Zero tangents, intervening higher-order
terms, arbitrary-pair boundary equivalence and growing-output
learning retain their separate hypotheses and questions.
